% introduction.tex -- EVM Executor Paper (BNR-2026-008)

\section{Introduction}

The construction of a zero-knowledge Ethereum Virtual Machine (zkEVM) Layer~2
requires solving a fundamental architectural problem: how to execute EVM bytecode
with full compatibility while simultaneously producing the structured execution
traces that a zero-knowledge prover requires for witness generation. This problem
sits at the intersection of systems engineering (EVM execution performance, state
management, tracing infrastructure) and applied cryptography (constraint system
design, hash function selection, proof generation).

Basis Network is an enterprise blockchain platform deployed on Avalanche,
targeting Latin American industrial and commercial clients. The platform operates
an application-specific L1 (Subnet-EVM, Chain~ID 43199) with zero-fee
permissioned transactions. The next phase of development---an enterprise zkEVM
L2---requires each client enterprise to operate its own execution chain with
ZK validity proofs verified on the Basis Network L1. This architecture demands
an execution engine that satisfies three properties simultaneously:

\begin{enumerate}
  \item \textbf{Full Cancun opcode compatibility}: Enterprise smart contracts
    compiled with Solidity 0.8.24+ targeting the Cancun EVM version must execute
    identically to mainnet Ethereum, including transient storage (EIP-1153),
    memory copy (EIP-5656), and blob operations (EIP-4844).
  \item \textbf{Structured trace generation}: Every state-modifying operation
    (storage reads/writes, balance changes, nonce updates, log emissions,
    contract creation) must be captured in a format consumable by the downstream
    Rust-based ZK prover via gRPC.
  \item \textbf{Enterprise-grade throughput}: The execution engine must sustain
    at least 1{,}000 transactions per second with trace generation enabled,
    matching the throughput requirements of industrial traceability and
    commercial ERP workloads.
\end{enumerate}

The research question addressed by this paper is:

\medskip
\noindent\textit{Can a minimal fork of go-ethereum execute EVM transactions with
custom state management, producing execution traces necessary for ZK witness
generation, while maintaining 100\% Cancun opcode compatibility and processing
1{,}000+ simple transactions per second?}
\medskip

This question is investigated through a systematic analysis of five production
L2 architectures, a complete mapping of Cancun EVM opcodes to their
zero-knowledge constraint costs, and the design of a custom ZK tracer
that leverages Geth's event-driven tracing hooks API.

The remainder of this paper is organized as follows. Section~II surveys related
work on zkEVM architectures and EVM fork strategies. Section~III presents the
system model and architectural constraints. Section~IV describes the
experimental methodology, including the opcode constraint mapping and tracer
design. Section~V details the experimental setup. Section~VI presents results
including projected throughput benchmarks and constraint cost analysis.
Section~VII discusses implications, invariants, and limitations. Section~VIII
concludes with recommendations for the Basis Network execution engine.
